% GENERATED FILE. Edit canonical lessons, then run npm run generate:books.
% canonical-source-hash: fnv1a64:b473193ec9cfdc61
% canonical-lessons: GE-C450-geldautomat, GE-C450-geheimzahl, GE-C450-kontoauszug, GE-C450-einloesen, GE-C450-auszahlen

\chapter{Geld und Konto}
\label{ch:ge-geld-und-konto}
\hlchaptermodality{450}

\begin{chapteropening}
\textbf{By the end of this chapter:} \emph{I can deal with money and an account: the cash machine, the PIN, the bank statement, redeeming a voucher and paying out.}

Complete the last lesson of chapter 450: I can deal with money and an account: the cash machine, the PIN, the bank statement, redeeming a voucher and paying out.
\end{chapteropening}

\section[der Geldautomat]{der Geldautomat — the cash machine, the ATM}

\label{lesson:GE-C450-geldautomat}



\begin{quote}
Before the new one: say the German for the card reader, then the German for the change.
\end{quote}

\subsection*{You'll want to know: der Geldautomat}
\textbf{der Geldautomat} — \textquotedblleft{}the cash machine, the ATM\textquotedblright{}.

\begin{quote}
Der Geldautomat ist direkt neben dem Supermarkt.
\end{quote}

\begin{quote}
\emph{The cash machine is right next to the supermarket.}
\end{quote}

Plural: \emph{die Geldautomaten}. In Austria: \emph{der Bankomat}.

\begin{cousinweb}
\emph{Geldautomat} is built from \emph{Geld} 'money' and \emph{Automat} 'machine' (EtymDB), which goes back to Greek \emph{autómatos} 'acting by itself'.
\end{cousinweb}

\begin{grammarlens}[title={What you've built}]
One more word for telling, explaining and describing at B1.
\end{grammarlens}

\subsection*{Your turn}
\begin{itemize}
\raggedright
  \item \emph{Say it:} \emph{der Geldautomat}
  \item \emph{Say it:} \emph{der Geldautomat}, once more
  \item \emph{Say it:} a sentence of your own with \emph{der Geldautomat}
  \item \emph{Recall:} say the German for the card reader, then the German for the change, then say \emph{der Geldautomat} again
\end{itemize}

\begin{tcolorbox}[breakable,colback=teal!4,colframe=teal!35!black,title={Before you move on}]
What does der Geldautomat mean? (\textquotedblleft{}The cash machine, the ATM\textquotedblright{}.) Say it once more.
\end{tcolorbox}
\section[die Geheimzahl]{die Geheimzahl — the PIN (number)}

\label{lesson:GE-C450-geheimzahl}



\begin{quote}
Before the new one: say the German for the change, then the German for the cash machine.
\end{quote}

\subsection*{You'll want to know: die Geheimzahl}
\textbf{die Geheimzahl} — \textquotedblleft{}the PIN (number)\textquotedblright{}.

\begin{quote}
Er hat seine Geheimzahl vergessen.
\end{quote}

\begin{quote}
\emph{He has forgotten his PIN.}
\end{quote}

Plural: \emph{die Geheimzahlen}. Also \emph{die PIN}.

\begin{cousinweb}
\emph{Geheimzahl} is built from \emph{geheim} 'secret', which is usually explained as first meaning 'belonging to the home', from \emph{Heim} 'home', and \emph{Zahl} 'number', inherited from Old High German \emph{zala} (EtymDB).
\end{cousinweb}

\begin{grammarlens}[title={What you've built}]
One more word for telling, explaining and describing at B1.
\end{grammarlens}

\subsection*{Your turn}
\begin{itemize}
\raggedright
  \item \emph{Say it:} \emph{die Geheimzahl}
  \item \emph{Say it:} \emph{die Geheimzahl}, once more
  \item \emph{Say it:} a sentence of your own with \emph{die Geheimzahl}
  \item \emph{Recall:} say the German for the change, then the German for the cash machine, then say \emph{die Geheimzahl} again
\end{itemize}

\begin{tcolorbox}[breakable,colback=teal!4,colframe=teal!35!black,title={Before you move on}]
What does die Geheimzahl mean? (\textquotedblleft{}The PIN (number)\textquotedblright{}.) Say it once more.
\end{tcolorbox}
\section[der Kontoauszug]{der Kontoauszug — the bank statement}

\label{lesson:GE-C450-kontoauszug}



\begin{quote}
Before the new one: say the German for the cash machine, then the German for the PIN.
\end{quote}

\subsection*{You'll want to know: der Kontoauszug}
\textbf{der Kontoauszug} — \textquotedblleft{}the bank statement\textquotedblright{}.

\begin{quote}
Auf dem Kontoauszug sieht sie alle Zahlungen.
\end{quote}

\begin{quote}
\emph{She can see all the payments on her bank statement.}
\end{quote}

Plural: \emph{die Kontoauszüge}.

\begin{cousinweb}
\emph{Kontoauszug} is built from \emph{Konto} 'account', borrowed from Italian \emph{conto} 'calculation' (EtymDB), which goes back to Latin \emph{computāre} 'to reckon', and \emph{Auszug} 'extract', from \emph{ausziehen} 'to pull out'.
\end{cousinweb}

\begin{grammarlens}[title={What you've built}]
One more word for telling, explaining and describing at B1.
\end{grammarlens}

\subsection*{Your turn}
\begin{itemize}
\raggedright
  \item \emph{Say it:} \emph{der Kontoauszug}
  \item \emph{Say it:} \emph{der Kontoauszug}, once more
  \item \emph{Say it:} a sentence of your own with \emph{der Kontoauszug}
  \item \emph{Recall:} say the German for the cash machine, then the German for the PIN, then say \emph{der Kontoauszug} again
\end{itemize}

\begin{tcolorbox}[breakable,colback=teal!4,colframe=teal!35!black,title={Before you move on}]
What does der Kontoauszug mean? (\textquotedblleft{}The bank statement\textquotedblright{}.) Say it once more.
\end{tcolorbox}
\section[einlösen]{einlösen — to redeem (a voucher), to cash (a cheque)}

\label{lesson:GE-C450-einloesen}



\begin{quote}
Before the new one: say the German for the PIN, then the German for the bank statement.
\end{quote}

\subsection*{You'll want to know: einlösen}
\textbf{einlösen} — \textquotedblleft{}to redeem (a voucher), to cash (a cheque)\textquotedblright{}.

\begin{quote}
An der Kasse löst sie einen Gutschein ein.
\end{quote}

\begin{quote}
\emph{At the till she redeems a voucher.}
\end{quote}

Separable: \emph{ich löse ein}, \emph{ich habe eingelöst}. \emph{Ein Versprechen einlösen} — to keep a promise.

\begin{cousinweb}
\emph{einlösen} is built from \emph{ein} 'in' and \emph{lösen} 'to release', inherited from Proto-Germanic \emph{lausijaną}, built on \emph{lausaz} 'loose' (EtymDB): to get back what a pledge held.
\end{cousinweb}

\begin{grammarlens}[title={What you've built}]
One more word for telling, explaining and describing at B1.
\end{grammarlens}

\subsection*{Your turn}
\begin{itemize}
\raggedright
  \item \emph{Say it:} \emph{einlösen}
  \item \emph{Say it:} \emph{einlösen}, once more
  \item \emph{Say it:} a sentence of your own with \emph{einlösen}
  \item \emph{Recall:} say the German for the PIN, then the German for the bank statement, then say \emph{einlösen} again
\end{itemize}

\begin{tcolorbox}[breakable,colback=teal!4,colframe=teal!35!black,title={Before you move on}]
What does einlösen mean? (\textquotedblleft{}To redeem (a voucher), to cash (a cheque)\textquotedblright{}.) Say it once more.
\end{tcolorbox}
\section[auszahlen]{auszahlen — to pay out}

\label{lesson:GE-C450-auszahlen}



\begin{quote}
Before the new one: say the German for the bank statement, then the German for to redeem.
\end{quote}

\subsection*{You'll want to know: auszahlen}
\textbf{auszahlen} — \textquotedblleft{}to pay out\textquotedblright{}.

\begin{quote}
Der Automat zahlt nur Scheine aus.
\end{quote}

\begin{quote}
\emph{The machine only pays out notes.}
\end{quote}

Separable: \emph{zahlt aus}, \emph{hat ausgezahlt}. \emph{Sich auszahlen} — to pay off, to be worth it.

\begin{cousinweb}
\emph{auszahlen} is built from \emph{aus} 'out' and \emph{zahlen} 'to pay', from \emph{Zahl} 'number', inherited from Old High German \emph{zala} (EtymDB).
\end{cousinweb}

\begin{grammarlens}[title={What you've built}]
One more word for telling, explaining and describing at B1.
\end{grammarlens}

\subsection*{Your turn}
\begin{itemize}
\raggedright
  \item \emph{Say it:} \emph{auszahlen}
  \item \emph{Say it:} \emph{auszahlen}, once more
  \item \emph{Say it:} a sentence of your own with \emph{auszahlen}
  \item \emph{Recall:} say the German for the bank statement, then the German for to redeem, then say \emph{auszahlen} again
\end{itemize}

\begin{tcolorbox}[breakable,colback=teal!4,colframe=teal!35!black,title={Before you move on}]
What does auszahlen mean? (\textquotedblleft{}To pay out\textquotedblright{}.) Say it once more.
\end{tcolorbox}
